\documentclass[11pt]{article}

\usepackage[margin=1in]{geometry}
\usepackage{amsmath,amssymb,amsthm,mathtools}
\usepackage{booktabs}
\usepackage{enumitem}
\usepackage[numbers,sort&compress]{natbib}
\usepackage{hyperref}
\usepackage{microtype}
\usepackage{xcolor}

\hypersetup{colorlinks=true,linkcolor=blue,citecolor=blue,urlcolor=blue}
\allowdisplaybreaks

\newtheorem{proposition}{Proposition}
\newtheorem{lemma}{Lemma}
\newtheorem{definition}{Definition}
\numberwithin{equation}{section}

\newcommand{\R}{\mathbb R}
\newcommand{\E}{\mathbb E}
\newcommand{\N}{\mathcal N}
\newcommand{\dd}{\,\mathrm d}
\newcommand{\tr}{\operatorname{tr}}
\newcommand{\diag}{\operatorname{diag}}
\newcommand{\vecop}{\operatorname{vec}}
\newcommand{\sgq}{\operatorname{SGQ}}
\newcommand{\chol}{\operatorname{chol}}
\newcommand{\ellhat}{\widehat\ell}
\newcommand{\bfone}{\mathbf 1}

\title{A Standalone Companion To Fixed Sparse-Grid Quadrature Filtering And Its Analytical Gradient}
\author{Chak Wong}
\date{2026-06-03}

\begin{document}
\maketitle

\begin{abstract}
This note develops fixed sparse-grid quadrature filtering as a standalone
high-dimensional filtering proposal.  The construction starts from the
sparse-grid quadrature nonlinear filter of Jia, Xin, and Cheng
\cite{jia2012sparsegrid}, expands the Gaussian projection and sparse-grid
formulas in BayesFilter notation, and then defines a fixed approximate
likelihood scalar with an analytical same-scalar gradient.  The scalar is a
deterministic Gaussian-projection likelihood: it is a candidate fixed target
for HMC experiments when branch diagnostics remain stable, but it is not the
exact nonlinear filtering likelihood and it does not store the full posterior
density.  The final section compares this proposal with the Zhao--Cui squared
tensor-train proposal \cite{zhao2024ttsequential}: fixed SGQF is transparent
and fixed-rule differentiable, while squared TT is a richer non-Gaussian
density approximation.
\end{abstract}

\tableofcontents

\section{What This Note Computes}
\label{sec:p31-computes}

The object computed here is deliberately narrow.  At time \(t\), before
observing \(y_t\), the filter carries one Gaussian approximation
\begin{equation}
  X_{t-1}\mid y_{1:t-1}
  \approx
  \N(m_{t-1},P_{t-1}).
  \label{eq:p31-carried-gaussian}
\end{equation}
It then uses a deterministic sparse-grid quadrature rule to approximate the
Gaussian moments needed by a Kalman-style projection update.  The output after
observing \(y_t\) is again only
\begin{equation}
  X_t\mid y_{1:t}
  \approx
  \N(m_t,P_t),
  \label{eq:p31-output-gaussian}
\end{equation}
together with the approximate Gaussian innovation log likelihood
\begin{equation}
  \ellhat_t
  =
  -\frac12
  \left\{
  \log\det S_t+v_t^\top S_t^{-1}v_t+n_y\log(2\pi)
  \right\}.
  \label{eq:p31-innovation-loglik-opening}
\end{equation}
The fixed-SGQF proposal is to use the sum
\begin{equation}
  \ellhat_T^{\rm FSGQ}(\theta)=\sum_{t=1}^T \ellhat_t^{\rm FSGQ}(\theta)
  \label{eq:p31-total-fixed-scalar-opening}
\end{equation}
as a deterministic approximate likelihood scalar and to differentiate that
same scalar analytically.

This note therefore computes an approximate likelihood and a Gaussian
projection recursion.  It does not compute the full nonlinear posterior law.
Sparse-grid exactness is exactness of a quadrature formula for selected
Gaussian-weighted polynomial integrands.  It is not a guarantee that the
posterior is Gaussian, unimodal, or accurately represented by \((m_t,P_t)\).

\paragraph{A scalar cell that keeps us honest.}
Let
\begin{equation}
  X\sim\N(0,P),\qquad Y=X^2+\varepsilon,\qquad
  \varepsilon\sim\N(0,R).
  \label{eq:p31-quadratic-cell}
\end{equation}
For \(y>0\) and small \(R\), the exact posterior density is proportional to
\begin{equation}
  \exp\!\left[-\frac{x^2}{2P}
  -\frac{(y-x^2)^2}{2R}\right],
  \label{eq:p31-quadratic-posterior}
\end{equation}
which can have two separated lobes near \(+\sqrt y\) and \(-\sqrt y\).  A
Gaussian projection filter can compute accurate moments of \(X^2\), but after
the update it still stores only one Gaussian.  This is not a bug in the
quadrature rule.  It is the modeling choice in \eqref{eq:p31-output-gaussian}.

The same limitation explains the role of fixed SGQF in the high-dimensional
program.  It is a transparent fixed approximate likelihood target with a
well-defined gradient.  It is not the non-Gaussian density approximation lane;
that role belongs to squared tensor-train filtering.

\section{State-Space Model And Exact Filtering Recursion}
\label{sec:p31-ssm}

Use the additive-noise state-space model
\begin{align}
  X_t &= f_\theta(X_{t-1})+\eta_t,
  &
  \eta_t&\sim\N(0,Q_\theta),
  \label{eq:p31-state-model}\\
  Y_t &= h_\theta(X_t)+\varepsilon_t,
  &
  \varepsilon_t&\sim\N(0,R_\theta),
  \label{eq:p31-observation-model}
\end{align}
with \(X_t\in\R^{n_x}\), \(Y_t\in\R^{n_y}\), and parameter \(\theta\in\R^p\).
The exact Bayesian prediction and update are
\begin{align}
  p_\theta(x_t\mid y_{1:t-1})
  &=
  \int p_\theta(x_t\mid x_{t-1})
       p_\theta(x_{t-1}\mid y_{1:t-1})\dd x_{t-1},
  \label{eq:p31-exact-prediction}\\
  p_\theta(x_t\mid y_{1:t})
  &=
  \frac{
  g_\theta(y_t\mid x_t)
  p_\theta(x_t\mid y_{1:t-1})}
  {\int g_\theta(y_t\mid x_t)
  p_\theta(x_t\mid y_{1:t-1})\dd x_t}.
  \label{eq:p31-exact-update}
\end{align}
These are the conceptual filtering equations used by Gaussian approximation
filters in Jia, Xin, and Cheng \cite[Section 2]{jia2012sparsegrid}.  In
nonlinear models, the integrals in \eqref{eq:p31-exact-prediction} and
\eqref{eq:p31-exact-update} usually do not close in finite-dimensional form.
Fixed SGQF replaces those exact densities by a Gaussian moment projection.

\section{Gaussian Projection From Moments}
\label{sec:p31-gaussian-projection}

Assume the predictive state is represented by
\begin{equation}
  X_t^-\sim\N(m_t^-,P_t^-).
  \label{eq:p31-predictive-gaussian}
\end{equation}
Let
\begin{equation}
  Z_t=h_\theta(X_t^-)
  \label{eq:p31-predicted-observation}
\end{equation}
be the noiseless predicted observation.  The Gaussian projection needs
\begin{align}
  \bar z_t &= \E[Z_t],
  \label{eq:p31-zbar}\\
  S_t &= R_\theta+\E[(Z_t-\bar z_t)(Z_t-\bar z_t)^\top],
  \label{eq:p31-S}\\
  C_{xz,t}&=\E[(X_t^- -m_t^-)(Z_t-\bar z_t)^\top].
  \label{eq:p31-Cxz}
\end{align}
Given these moments, define
\begin{align}
  v_t&=y_t-\bar z_t,
  &
  K_t&=C_{xz,t}S_t^{-1},
  \label{eq:p31-v-K}\\
  m_t&=m_t^-+K_tv_t,
  &
  P_t&=P_t^- - K_tS_tK_t^\top.
  \label{eq:p31-projection-update}
\end{align}

\begin{proposition}[Affine projection]
\label{prop:p31-affine-projection}
Among estimators of \(X_t^-\) of the form \(a+B Y_t\), the minimum mean-square
estimator under the moments in
\eqref{eq:p31-zbar}--\eqref{eq:p31-Cxz} uses \(B=C_{xz,t}S_t^{-1}\).  Storing
the resulting mean and covariance as a Gaussian gives
\eqref{eq:p31-projection-update}.
\end{proposition}

\begin{proof}
Center \(\widetilde X=X_t^- -m_t^-\) and
\(\widetilde Y=Y_t-\bar z_t\).  The squared-error objective for
\(m_t^-+B\widetilde Y\) is
\begin{equation}
  J(B)=\E\|\widetilde X-B\widetilde Y\|^2
  =
  \tr(P_t^-)-2\tr(BC_{xz,t}^\top)+\tr(BS_tB^\top).
  \label{eq:p31-affine-objective}
\end{equation}
Differentiating gives \(\nabla_BJ=-2C_{xz,t}+2BS_t\).  If \(S_t\) is
nonsingular, the stationary point is \(B=C_{xz,t}S_t^{-1}\).  Substitution gives
the residual covariance \(P_t^- - C_{xz,t}S_t^{-1}C_{xz,t}^\top\), which equals
\eqref{eq:p31-projection-update}.
\end{proof}

The approximation has two layers.  First, the moments
\eqref{eq:p31-zbar}--\eqref{eq:p31-Cxz} may be approximated by quadrature.
Second, the resulting posterior information is compressed to one Gaussian.

\section{One-Dimensional Gaussian Quadrature And Tensor Products}
\label{sec:p31-ghq}

Let \(\xi\sim\N(0,1)\).  A one-dimensional standard-normal quadrature rule at
level \(\ell\) is a set of nodes and weights
\begin{equation}
  I_\ell[\psi]
  =
  \sum_{a=1}^{s_\ell}\omega_{\ell,a}\psi(\zeta_{\ell,a})
  \approx
  \int_{\R}\psi(\xi)\phi(\xi)\dd \xi.
  \label{eq:p31-univariate-rule}
\end{equation}
For a multi-index \(\mathbf i=(i_1,\ldots,i_b)\), the corresponding tensor rule
is
\begin{equation}
  I_{\mathbf i}[F]
  =
  (I_{i_1}\otimes\cdots\otimes I_{i_b})[F].
  \label{eq:p31-tensor-rule}
\end{equation}
Its nodes and weights are
\begin{align}
  \xi^{\mathbf i,\alpha}
  &=
  (\zeta_{i_1,\alpha_1},\ldots,\zeta_{i_b,\alpha_b}),
  \label{eq:p31-tensor-node}\\
  w^{\mathbf i,\alpha}
  &=
  \prod_{j=1}^b\omega_{i_j,\alpha_j},
  \qquad
  \alpha_j\in\{1,\ldots,s_{i_j}\}.
  \label{eq:p31-tensor-weight}
\end{align}
A full tensor-product rule is simple but expensive: if every coordinate uses
\(s\) nodes, then \(M=s^b\).  Jia, Xin, and Cheng introduce the sparse-grid
combination to reduce this point count for fixed accuracy level
\cite[Section 3]{jia2012sparsegrid}.

\section{Sparse-Grid Rule Reconstructed}
\label{sec:p31-sgq-reconstruction}

For dimension \(b\) and level \(L\), define the level band
\begin{equation}
  \mathcal A_{b,L}
  =
  \{\mathbf i\in\mathbb N^b: L\le |\mathbf i|\le L+b-1\},
  \qquad |\mathbf i|=\sum_{j=1}^b i_j.
  \label{eq:p31-level-band}
\end{equation}
For each \(\mathbf i\in\mathcal A_{b,L}\), set
\begin{equation}
  c_{\mathbf i}
  =
  (-1)^{L+b-1-|\mathbf i|}
  \binom{b-1}{L+b-1-|\mathbf i|}.
  \label{eq:p31-smolyak-coeff}
\end{equation}
The sparse-grid quadrature approximation is
\begin{equation}
  A_{b,L}[F]
  =
  \sum_{\mathbf i\in\mathcal A_{b,L}}
  c_{\mathbf i}I_{\mathbf i}[F].
  \label{eq:p31-smolyak-rule}
\end{equation}
This is the BayesFilter notation for the Jia--Xin--Cheng sparse-grid formula
\cite[Eq. 26--29]{jia2012sparsegrid}.  Their notation groups indices by
\(|\mathbf i|=b+q\); the band \(q=L-b,\ldots,L-1\) is equivalent to
\eqref{eq:p31-level-band}.

\paragraph{From formula to cloud.}
Equation \eqref{eq:p31-smolyak-rule} is not yet an implementation object.  The
same node can appear in several tensor rules.  The fixed cloud is the
accumulated dictionary
\begin{equation}
  D[\xi]
  =
  \sum_{\substack{\mathbf i,\alpha:
  \xi^{\mathbf i,\alpha}=\xi}}
  c_{\mathbf i}w^{\mathbf i,\alpha}.
  \label{eq:p31-node-dictionary}
\end{equation}
After merging duplicate standardized nodes, enumerate the nonzero entries as
\begin{equation}
  \mathcal C_{b,L}
  =
  \{(\xi^{(r)},w_r)\}_{r=1}^{M_{b,L}}.
  \label{eq:p31-fixed-cloud}
\end{equation}
Jia, Xin, and Cheng describe this point-and-weight generation in Algorithm 1
\cite[Algorithm 1]{jia2012sparsegrid}.  For the fixed BayesFilter scalar, the
cloud \(\mathcal C_{b,L}\), the one-dimensional rules, the merge tolerance, and
the accumulated signed weights are part of the mathematical target.

\paragraph{Source exactness and point count.}
The source exactness statement is a quadrature statement: under the univariate
conditions in \citet[Theorem 3.1]{jia2012sparsegrid}, the sparse-grid rule
inherits a polynomial exactness level.  This does not imply exact nonlinear
posterior filtering.  The point-count statements in
\citet[Proposition 3.1--3.2]{jia2012sparsegrid} are fixed-level statements.
They support the claim that sparse grids can grow polynomially in dimension
for fixed level, not the claim that a given fixed level is accurate for every
high-dimensional nonlinear posterior.

\paragraph{UKF relation.}
Jia, Xin, and Cheng prove that an unscented-transform point set is a member of
the SGQF family at level 2 under their parameterization
\cite[Theorem 3.2]{jia2012sparsegrid}.  In this note, the relation is used only
as orientation: fixed SGQF generalizes the idea of deterministic Gaussian
moment points, but the selected BayesFilter target is the declared fixed cloud
\eqref{eq:p31-fixed-cloud}.

\section{A Toy Fixed Grid With Duplicate Merging}
\label{sec:p31-toy-grid}

This section makes the sparse-grid cloud concrete.  Take \(b=2\) and \(L=2\).
Use
\begin{equation}
  I_1:\quad (0,1),
  \qquad
  I_2:\quad
  \left\{(-\sqrt3,1/6),(0,2/3),(\sqrt3,1/6)\right\},
  \label{eq:p31-toy-univariate}
\end{equation}
which is the familiar three-point standard-normal rule at level 2.  The band
\(\mathcal A_{2,2}\) contains
\begin{equation}
  (1,1),\qquad (1,2),\qquad (2,1),
  \label{eq:p31-toy-indices}
\end{equation}
with coefficients
\begin{equation}
  c_{(1,1)}=-1,\qquad c_{(1,2)}=1,\qquad c_{(2,1)}=1.
  \label{eq:p31-toy-coefficients}
\end{equation}
Before merging, the node \((0,0)\) appears three times.  Its accumulated weight
is
\begin{equation}
  -1+\frac23+\frac23=\frac13.
  \label{eq:p31-toy-origin-weight}
\end{equation}
The four axis nodes each receive weight \(1/6\).  The merged cloud is
\begin{center}
\begin{tabular}{ccc}
\toprule
node \(\xi\) & accumulated weight & source components \\
\midrule
\((0,0)\) & \(1/3\) & \((1,1),(1,2),(2,1)\) \\
\((-\sqrt3,0)\) & \(1/6\) & \((2,1)\) \\
\((\sqrt3,0)\) & \(1/6\) & \((2,1)\) \\
\((0,-\sqrt3)\) & \(1/6\) & \((1,2)\) \\
\((0,\sqrt3)\) & \(1/6\) & \((1,2)\) \\
\bottomrule
\end{tabular}
\end{center}
The weights sum to one:
\begin{equation}
  \frac13+4\cdot\frac16=1.
  \label{eq:p31-toy-weight-total}
\end{equation}
This toy cloud illustrates two implementation rules that are easy to get
wrong.  First, duplicate nodes must be accumulated before the cloud is used.
Second, signed sparse-grid coefficients are applied before accumulation; a
future level or a different univariate family can produce negative accumulated
weights, so covariance matrices require explicit diagnostics.

\section{FixedSGQF Filtering Value Path}
\label{sec:p31-value-path}

Fix a standardized sparse-grid cloud
\begin{equation}
  \mathcal C=\{(\xi^{(r)},w_r)\}_{r=1}^M,
  \qquad
  \xi^{(r)}\in\R^{n_x},
  \quad
  w_r\in\R.
  \label{eq:p31-cloud-for-filter}
\end{equation}
The cloud is frozen before likelihood evaluation.  The value path has two
quadrature stages at each time.

\subsection{Prediction Stage}

Let \(P_{t-1}=C_{t-1}C_{t-1}^\top\) on the declared square-root branch.  Place
the previous Gaussian cloud:
\begin{equation}
  \chi_{t-1}^{(r)}
  =
  m_{t-1}+C_{t-1}\xi^{(r)}.
  \label{eq:p31-pred-placed-points}
\end{equation}
Push points through the transition:
\begin{equation}
  a_t^{(r)}=f_\theta(\chi_{t-1}^{(r)}).
  \label{eq:p31-transition-points}
\end{equation}
The predicted Gaussian is
\begin{align}
  m_t^-&=\sum_{r=1}^Mw_ra_t^{(r)},
  \label{eq:p31-pred-mean}\\
  P_t^-&=Q_\theta+\sum_{r=1}^Mw_r
  (a_t^{(r)}-m_t^-)(a_t^{(r)}-m_t^-)^\top.
  \label{eq:p31-pred-cov}
\end{align}
After \eqref{eq:p31-pred-cov}, the implementation replaces \(P_t^-\) by
\((P_t^-+P_t^{-\top})/2\) and checks the declared positive-definite branch.
For the differentiable lane in this note, there is no adaptive jitter, floor,
clipping, pivot, or eigenvalue repair inside the scalar.  If the check fails,
the value path returns a veto rather than quietly changing the scalar.

\subsection{Observation Stage}

Factor \(P_t^-=C_t^-(C_t^-)^\top\) on the declared branch and place the
predictive cloud:
\begin{equation}
  \chi_t^{(r)}
  =
  m_t^-+C_t^-\xi^{(r)}.
  \label{eq:p31-obs-placed-points}
\end{equation}
Evaluate the observation map:
\begin{equation}
  z_t^{(r)}=h_\theta(\chi_t^{(r)}).
  \label{eq:p31-obs-values}
\end{equation}
Then compute
\begin{align}
  \bar z_t&=\sum_{r=1}^Mw_rz_t^{(r)},
  \label{eq:p31-obs-mean}\\
  \Delta_t^{(r)}&=z_t^{(r)}-\bar z_t,
  \label{eq:p31-obs-centered}\\
  S_t&=R_\theta+\sum_{r=1}^Mw_r
  \Delta_t^{(r)}\Delta_t^{(r)\top},
  \label{eq:p31-obs-cov}\\
  C_{xz,t}&=\sum_{r=1}^Mw_r
  (\chi_t^{(r)}-m_t^-)\Delta_t^{(r)\top}.
  \label{eq:p31-cross-cov}
\end{align}
The innovation and log likelihood are
\begin{align}
  v_t&=y_t-\bar z_t,
  \label{eq:p31-innovation}\\
  \ellhat_t^{\rm FSGQ}
  &=
  -\frac12
  \left\{\log\det S_t+v_t^\top S_t^{-1}v_t+n_y\log(2\pi)\right\}.
  \label{eq:p31-fsgq-loglik}
\end{align}
If \(S_t\) is not positive definite on the declared branch, the value path
returns a veto.  If it passes, the Gaussian projection update is
\begin{align}
  K_t&=C_{xz,t}S_t^{-1},
  \label{eq:p31-kalman-gain}\\
  m_t&=m_t^-+K_tv_t,
  \label{eq:p31-filter-mean}\\
  P_t&=P_t^- - K_tS_tK_t^\top.
  \label{eq:p31-filter-cov}
\end{align}

\section{The Saved Scalar And Same-Scalar Contract}
\label{sec:p31-same-scalar}

The derivative in this note is meaningful only for the scalar actually
evaluated.  A fixed-SGQF branch is the tuple
\begin{equation}
\begin{aligned}
  \mathcal B
  =
  (&y_{1:T},\ \text{observation preprocessing},\
  m_0(\theta),P_0(\theta),\dot m_0,\dot P_0,\\
  &\mathcal C,\ \text{one-dimensional rules},\
  \text{duplicate-merge tolerance},\
  C(\cdot),\\
  &\text{symmetrize-then-veto rule},\
  f_\theta,h_\theta,Q_\theta,R_\theta).
\end{aligned}
  \label{eq:p31-branch-tuple}
\end{equation}
The fixed scalar is
\begin{equation}
  \ellhat_T^{\rm FSGQ}(\theta;\mathcal B)
  =
  \sum_{t=1}^T \ellhat_t^{\rm FSGQ}(\theta;\mathcal B),
  \label{eq:p31-fixed-scalar}
\end{equation}
where every term is computed by
\eqref{eq:p31-pred-placed-points}--\eqref{eq:p31-filter-cov}.  The same
\(\mathcal B\) must be used for values, gradients, and finite differences.
Changing the grid level, active index set, duplicate-merge tolerance,
Cholesky pivot, observation preprocessing, initial-condition policy, eigenvalue
floor, covariance clipping, or point pruning changes the scalar.  The
derivative formulas below cover the open branch where the symmetrized
covariances pass the Cholesky and positive-definite checks without repair.  If
a jitter, floor, clipping, or pivoting map is desired, it must be declared as a
different scalar with its own derivative or stop rule.

\section{Analytical Gradient Of The Fixed Scalar}
\label{sec:p31-gradient}

Differentiate with respect to parameter component \(\theta_i\).  A dot denotes
\(\partial_i\), for example \(\dot m_t=\partial_i m_t\).  The data \(y_t\) are
recorded and do not depend on \(\theta\).

\subsection{Square-Root Branch}

For the default unpivoted Cholesky branch, \(P=LL^\top\) with positive diagonal
and \(C(P)=L\).  Given symmetric \(\dot P\), set
\begin{equation}
  A=L^{-1}\dot P L^{-\top}.
  \label{eq:p31-chol-A}
\end{equation}
Define the lower-triangular operator
\begin{equation}
  \Phi(A)_{jk}
  =
  \begin{cases}
  A_{jk}, & j>k,\\
  A_{jj}/2, & j=k,\\
  0, & j<k.
  \end{cases}
  \label{eq:p31-chol-Phi}
\end{equation}
Then
\begin{equation}
  \dot L=L\Phi(A).
  \label{eq:p31-chol-derivative}
\end{equation}
Indeed, \(\Phi(A)+\Phi(A)^\top=A\), so
\begin{equation}
  \dot LL^\top+L\dot L^\top
  =
  L\{\Phi(A)+\Phi(A)^\top\}L^\top
  =
  \dot P.
  \label{eq:p31-chol-check}
\end{equation}
If a different square-root map is used, the implementation must provide
\(\dot C=DC(P)[\dot P]\) for the same branch.

\subsection{Prediction Sensitivities}

From \eqref{eq:p31-pred-placed-points},
\begin{equation}
  \dot\chi_{t-1}^{(r)}
  =
  \dot m_{t-1}+\dot C_{t-1}\xi^{(r)}.
  \label{eq:p31-dot-pred-place}
\end{equation}
From \eqref{eq:p31-transition-points},
\begin{equation}
  \dot a_t^{(r)}
  =
  D_xf_\theta(\chi_{t-1}^{(r)})\dot\chi_{t-1}^{(r)}
  +\partial_i f_\theta(\chi_{t-1}^{(r)}).
  \label{eq:p31-dot-transition}
\end{equation}
The predicted mean derivative is
\begin{equation}
  \dot m_t^-=\sum_{r=1}^M w_r\dot a_t^{(r)}.
  \label{eq:p31-dot-pred-mean}
\end{equation}
Let
\begin{equation}
  d_t^{(r)}=a_t^{(r)}-m_t^-,
  \qquad
  \dot d_t^{(r)}=\dot a_t^{(r)}-\dot m_t^-.
  \label{eq:p31-pred-centered}
\end{equation}
Then
\begin{equation}
  \dot P_t^-
  =
  \dot Q_\theta
  +
  \sum_{r=1}^M w_r
  \left\{
  \dot d_t^{(r)}d_t^{(r)\top}
  +d_t^{(r)}\dot d_t^{(r)\top}
  \right\}.
  \label{eq:p31-dot-pred-cov}
\end{equation}

\subsection{Observation Sensitivities}

Differentiate \eqref{eq:p31-obs-placed-points}:
\begin{equation}
  \dot\chi_t^{(r)}
  =
  \dot m_t^-+\dot C_t^-\xi^{(r)}.
  \label{eq:p31-dot-obs-place}
\end{equation}
Then
\begin{equation}
  \dot z_t^{(r)}
  =
  D_xh_\theta(\chi_t^{(r)})\dot\chi_t^{(r)}
  +\partial_i h_\theta(\chi_t^{(r)}).
  \label{eq:p31-dot-obs-value}
\end{equation}
The observation mean, residual, and centered residual derivatives are
\begin{align}
  \dot{\bar z}_t&=\sum_{r=1}^M w_r\dot z_t^{(r)},
  \label{eq:p31-dot-zbar}\\
  \dot v_t&=-\dot{\bar z}_t,
  \label{eq:p31-dot-v}\\
  \dot\Delta_t^{(r)}&=\dot z_t^{(r)}-\dot{\bar z}_t.
  \label{eq:p31-dot-Delta}
\end{align}
Differentiate \(S_t\) and \(C_{xz,t}\):
\begin{align}
  \dot S_t
  &=
  \dot R_\theta
  +
  \sum_{r=1}^M w_r
  \left\{
  \dot\Delta_t^{(r)}\Delta_t^{(r)\top}
  +\Delta_t^{(r)}\dot\Delta_t^{(r)\top}
  \right\},
  \label{eq:p31-dot-S}\\
  \dot C_{xz,t}
  &=
  \sum_{r=1}^M w_r
  \left\{
  (\dot\chi_t^{(r)}-\dot m_t^-)\Delta_t^{(r)\top}
  +(\chi_t^{(r)}-m_t^-)\dot\Delta_t^{(r)\top}
  \right\}.
  \label{eq:p31-dot-Cxz}
\end{align}

\subsection{Innovation Score}

\begin{proposition}[Fixed Gaussian innovation score]
\label{prop:p31-innovation-score}
Assume \(S_t\) is positive definite and the active branch is differentiable.
Let \(u_t\) solve
\begin{equation}
  S_tu_t=v_t.
  \label{eq:p31-score-solve}
\end{equation}
Then the derivative of \eqref{eq:p31-fsgq-loglik} is
\begin{equation}
  \dot\ellhat_t^{\rm FSGQ}
  =
  -\frac12
  \left[
  \tr(S_t^{-1}\dot S_t)
  +2\dot v_t^\top u_t
  -u_t^\top\dot S_tu_t
  \right].
  \label{eq:p31-score}
\end{equation}
\end{proposition}

\begin{proof}
Use
\begin{equation}
  \dd\log\det S=\tr(S^{-1}\dd S)
  \label{eq:p31-d-logdet}
\end{equation}
and
\begin{equation}
  \dd S^{-1}=-S^{-1}(\dd S)S^{-1}.
  \label{eq:p31-d-inverse}
\end{equation}
Therefore, with \(u=S^{-1}v\),
\begin{equation}
  \dd(v^\top S^{-1}v)
  =
  2(\dd v)^\top u-u^\top(\dd S)u.
  \label{eq:p31-d-quadratic}
\end{equation}
Substitute \(\dd S=\dot S_t\) and \(\dd v=\dot v_t\) into
\eqref{eq:p31-innovation-loglik-opening}.
\end{proof}

\subsection{Posterior Sensitivity Propagation}

The gradient must also propagate \(\dot m_t\) and \(\dot P_t\) into the next
time step.  Differentiate \(K_t=C_{xz,t}S_t^{-1}\):
\begin{equation}
  \dot K_t
  =
  \dot C_{xz,t}S_t^{-1}
  -K_t\dot S_tS_t^{-1}.
  \label{eq:p31-dot-K}
\end{equation}
Then
\begin{align}
  \dot m_t
  &=
  \dot m_t^-+\dot K_tv_t+K_t\dot v_t,
  \label{eq:p31-dot-m}\\
  \dot P_t
  &=
  \dot P_t^-
  -\dot K_tS_tK_t^\top
  -K_t\dot S_tK_t^\top
  -K_tS_t\dot K_t^\top.
  \label{eq:p31-dot-P}
\end{align}
Equations \eqref{eq:p31-dot-pred-place}--\eqref{eq:p31-dot-P} define the
value-and-gradient recursion for the fixed scalar
\eqref{eq:p31-fixed-scalar}.

\section{One Boxed Mathematical Algorithm}
\label{sec:p31-boxed-algorithm}

\fbox{\begin{minipage}{0.94\linewidth}
\textbf{FixedSGQF value and gradient for one saved branch.}

\begin{enumerate}[leftmargin=0.22in]
\item Choose dimension \(b=n_x\), level \(L\), one-dimensional rules
\(I_\ell\), and merge tolerance.  Build and save the cloud
\(\mathcal C=\{(\xi^{(r)},w_r)\}_{r=1}^M\) by
\eqref{eq:p31-level-band}--\eqref{eq:p31-fixed-cloud}.
\item Save \(y_{1:T}\), observation preprocessing, \(m_0(\theta),P_0(\theta)\),
their sensitivity policy, the factor map \(C(P)\), the symmetrize-then-veto
rule, and model derivative routines.
\item Initialize \(m_0,P_0\) and \(\dot m_0,\dot P_0\) for each parameter
component.
\item For \(t=1,\ldots,T\), run the prediction value equations
\eqref{eq:p31-pred-placed-points}--\eqref{eq:p31-pred-cov} and prediction
sensitivity equations \eqref{eq:p31-dot-pred-place}--\eqref{eq:p31-dot-pred-cov}.
\item If \(P_t^-\) violates the saved branch or veto policy, stop and report
the branch failure.
\item Run the observation value equations
\eqref{eq:p31-obs-placed-points}--\eqref{eq:p31-fsgq-loglik} and observation
sensitivity equations \eqref{eq:p31-dot-obs-place}--\eqref{eq:p31-dot-S}.
\item If \(S_t\) violates the saved positive-definite branch, stop and report
the branch failure.
\item Add the score \eqref{eq:p31-score} to the running gradient and add
\(\ellhat_t^{\rm FSGQ}\) to the running scalar.
\item Update \(m_t,P_t,\dot m_t,\dot P_t\) using
\eqref{eq:p31-kalman-gain}--\eqref{eq:p31-filter-cov} and
\eqref{eq:p31-dot-K}--\eqref{eq:p31-dot-P}.
\end{enumerate}
\end{minipage}}

\section{Implementation Contract}
\label{sec:p31-implementation-contract}

\begin{center}
\small
\begin{tabular}{p{0.18\linewidth}p{0.30\linewidth}p{0.16\linewidth}p{0.24\linewidth}}
\toprule
object & definition & shape & reuse \\
\midrule
\(\xi^{(r)}\) & standardized sparse-grid node & \(n_x\) & value and gradient \\
\(w_r\) & accumulated signed weight after duplicate merge & scalar & value and gradient \\
\(C_{t-1},C_t^-\) & saved covariance factors on declared branch & \(n_x\times n_x\) & value, gradient, diagnostics \\
\(\chi_{t-1}^{(r)}\) & previous placed point & \(n_x\) & transition value and derivative \\
\(a_t^{(r)}\) & transition image & \(n_x\) & prediction moments \\
\(\chi_t^{(r)}\) & predictive placed point & \(n_x\) & observation value and derivative \\
\(z_t^{(r)}\) & observation image & \(n_y\) & observation moments \\
\(\bar z_t,S_t,C_{xz,t}\) & Gaussian projection moments & \(n_y\), \(n_y\times n_y\), \(n_x\times n_y\) & likelihood and update \\
\(v_t,K_t\) & innovation and gain & \(n_y\), \(n_x\times n_y\) & likelihood, update, gradient \\
\(\dot \chi,\dot a,\dot z\) & point sensitivities & matching value objects & gradient only \\
\(\dot m,\dot P,\dot S,\dot C_{xz}\) & moment and filter sensitivities & matching value objects & gradient and next step \\
data record & \(y_{1:T}\), preprocessing, missing-data policy if any & metadata & value, gradient, finite difference \\
initial law & \(m_0(\theta),P_0(\theta),\dot m_0,\dot P_0\) & \(n_x,n_x\times n_x\) & value and gradient start \\
branch record & cloud, merge tolerance, factor branch, symmetrize-then-veto rule & metadata & same-scalar finite difference \\
\bottomrule
\end{tabular}
\end{center}

\section{Finite-Difference Same-Scalar Check}
\label{sec:p31-fd}

The finite-difference check must call the same value routine with the same
saved branch \(\mathcal B\).  For parameter component \(i\), use
\begin{equation}
  D_i(h)
  =
  \frac{
  \ellhat_T^{\rm FSGQ}(\theta+he_i;\mathcal B)
  -
  \ellhat_T^{\rm FSGQ}(\theta-he_i;\mathcal B)}
  {2h}.
  \label{eq:p31-central-diff}
\end{equation}
Use the step ladder
\begin{equation}
  h\in\{10^{-1},10^{-2},10^{-3},10^{-4},10^{-5}\}\,
  \max\{1,|\theta_i|\}.
  \label{eq:p31-fd-ladder}
\end{equation}
For every \(h\), the \(+\!h\) and \(-\!h\) evaluations must reuse the same
cloud, preprocessing, initial-condition policy, factor branch, merge tolerance,
and symmetrize-then-veto rule.  If either side hits a different branch or
returns a veto, that row is marked branch-invalid and is not interpreted as a
derivative check.  For valid rows, compare \(D_i(h)\) against the analytical
score \(G_i\) by
\begin{equation}
  e_i(h)=
  \frac{|D_i(h)-G_i|}
  {\max\{1,|D_i(h)|,|G_i|\}}.
  \label{eq:p31-fd-relative-error}
\end{equation}
The diagnostic passes when at least two consecutive valid rows have
\(e_i(h)\le 10^{-5}\) and the smaller-step row is no worse than ten times the
larger-step row.  This is a numerical diagnostic, not a theorem: failure means
one of three things is likely.  The derivative formula is wrong, the value and
gradient paths are not using the same scalar, or the branch is too nonsmooth or
ill-conditioned at that point.

\paragraph{A scalar numerical trace.}
Let \(X\sim\N(0,1)\), \(h_\theta(x)=\theta x\), \(R=1\), and \(y=2\).  Any
fixed Gaussian rule with exact second moment gives
\begin{equation}
  \bar z=0,\qquad S(\theta)=1+\theta^2,\qquad
  \ellhat(\theta)
  =
  -\frac12\left\{\log(1+\theta^2)+\frac{4}{1+\theta^2}+\log(2\pi)\right\}.
  \label{eq:p31-fd-toy-likelihood}
\end{equation}
At \(\theta=1\),
\begin{equation}
  G
  =
  -\frac12
  \left\{
  \frac{2}{2}
  -
  \frac{4\cdot2}{2^2}
  \right\}
  =
  0.5.
  \label{eq:p31-fd-toy-gradient}
\end{equation}
Central differences give:
\begin{center}
\begin{tabular}{rrrr}
\toprule
\(h\) & \(D(h)\) & \(G\) & \(|D(h)-G|\) \\
\midrule
\(10^{-1}\) & 0.5008108250 & 0.5 & \(8.10825\times10^{-4}\) \\
\(10^{-2}\) & 0.5000083311 & 0.5 & \(8.33108\times10^{-6}\) \\
\(10^{-3}\) & 0.5000000833 & 0.5 & \(8.33331\times10^{-8}\) \\
\(10^{-4}\) & 0.5000000008 & 0.5 & \(8.33722\times10^{-10}\) \\
\bottomrule
\end{tabular}
\end{center}

\paragraph{A cloud-sensitive nonlinear trace.}
The previous trace tests the innovation-score algebra.  To exercise the fixed
cloud, use the two-dimensional level-2 cloud in Section~\ref{sec:p31-toy-grid}
and the nonlinear observation
\begin{equation}
  h_\theta(x_1,x_2)=\theta(x_1^2+x_2^2),
  \qquad
  X\sim\N(0,I_2),\qquad R=1,\qquad y=3.
  \label{eq:p31-cloud-sensitive-model}
\end{equation}
For the merged cloud in Section~\ref{sec:p31-toy-grid},
\begin{equation}
  \sum_r w_r(x_{1,r}^2+x_{2,r}^2)=2,
  \qquad
  \sum_r w_r(x_{1,r}^2+x_{2,r}^2)^2=6.
  \label{eq:p31-cloud-sensitive-moments}
\end{equation}
Therefore
\begin{equation}
  \bar z(\theta)=2\theta,\qquad
  S(\theta)=1+2\theta^2,\qquad
  v(\theta)=3-2\theta,
  \label{eq:p31-cloud-sensitive-likelihood-objects}
\end{equation}
if the parameter scales the whole nonlinear observation linearly.  To exercise
the parameter through placed nonlinear values more sharply, set instead
\begin{equation}
  h_\theta(x_1,x_2)=\theta^2(x_1^2+x_2^2),
  \label{eq:p31-cloud-sensitive-square-param}
\end{equation}
which gives
\begin{equation}
  \bar z(\theta)=2\theta^2,\qquad
  S(\theta)=1+2\theta^4,\qquad
  v(\theta)=3-2\theta^2.
  \label{eq:p31-cloud-sensitive-square-objects}
\end{equation}
At \(\theta=1\), the score from \eqref{eq:p31-score} is
\begin{equation}
  G
  =
  -\frac12
  \left[
    \frac{8}{3}
    +2(-4)\frac{1}{3}
    -\frac{1}{3^2}8
  \right]
  =
  \frac49.
  \label{eq:p31-cloud-sensitive-score}
\end{equation}
Central differences for the saved cloud scalar give:
\begin{center}
\begin{tabular}{rrrr}
\toprule
\(h\) & \(D(h)\) & \(G\) & \(|D(h)-G|\) \\
\midrule
\(10^{-1}\) & 0.5200303321 & 0.4444444444 & \(7.55859\times10^{-2}\) \\
\(10^{-2}\) & 0.4452146703 & 0.4444444444 & \(7.70226\times10^{-4}\) \\
\(10^{-3}\) & 0.4444521481 & 0.4444444444 & \(7.70369\times10^{-6}\) \\
\(10^{-4}\) & 0.4444445215 & 0.4444444444 & \(7.70368\times10^{-8}\) \\
\bottomrule
\end{tabular}
\end{center}
This example depends on the merged sparse-grid fourth moment in
\eqref{eq:p31-cloud-sensitive-moments}.  A duplicate-merge bug or a different
saved cloud changes \(S(\theta)\), \(G\), and the finite-difference table.

\section{Diagnostics, Accuracy, Memory, And Performance Tests}
\label{sec:p31-validation}

Fixed SGQF should be tested as an approximate deterministic likelihood, not as
an exact posterior method.  The minimum validation suite is:

\begin{enumerate}[leftmargin=0.24in]
\item \textbf{Construction checks.}  Verify duplicate-node accumulation,
weight totals, signed-weight counts, and deterministic cloud reproduction for
small \((b,L)\) cases such as Section~\ref{sec:p31-toy-grid}.
\item \textbf{Linear-Gaussian recovery.}  For affine \(f_\theta,h_\theta\),
the Gaussian projection recursion should match the Kalman filter up to
quadrature and floating-point tolerance.
\item \textbf{Quadratic scalar cell.}  Use \eqref{eq:p31-quadratic-cell} to
show both the moment accuracy and the posterior-shape limitation.
\item \textbf{Same-scalar finite differences.}  Run
\eqref{eq:p31-central-diff} for representative parameters and step sizes.
\item \textbf{Signed-weight stability.}  Record the minimum eigenvalues of
symmetrized \(P_t^-\), \(S_t\), and \(P_t\), plus any branch veto.
\item \textbf{Memory and runtime scaling.}  Record \(M_{b,L}\), model
evaluations, wall time, and peak memory as \(b,L,T\) change.
\item \textbf{Accuracy ladder.}  Compare fixed SGQF against tensor-product GHQ
in small dimension, Kalman truth in linear-Gaussian models, and particle or TT
diagnostics in nonlinear models where those references are available.
\end{enumerate}

Passing these checks supports the statement that the implementation evaluates
and differentiates the declared approximate scalar.  It does not prove exact
posterior accuracy.

\section{Adaptive Sparse Grids As Grid Design}
\label{sec:p31-adaptive}

Singh et al. \cite{singh2018adaptive} choose sparse-grid indices adaptively by
using admissible index sets, local error indicators, active and old index
sets, a global error estimate, and a tolerance.  That is useful for designing a
cloud before inference.  For HMC, however, the live target cannot silently
change its grid as \(\theta\) moves.  The safe interpretation is
\begin{equation}
  \text{adaptive pilot chooses a cloud}
  \quad\Longrightarrow\quad
  \text{freeze the cloud}
  \quad\Longrightarrow\quad
  \text{run FixedSGQF}.
  \label{eq:p31-adaptive-to-fixed}
\end{equation}
If the adaptive index set is rerun inside the likelihood, the value path is a
piecewise-defined algorithm with discrete changes.  The formulas in
Section~\ref{sec:p31-gradient} do not differentiate those changes.

\section{Relation To Zhao--Cui Squared TT}
\label{sec:p31-comparison}

Fixed SGQF and Zhao--Cui squared TT solve different parts of the high-dimensional
filtering problem.

\begin{center}
\small
\begin{tabular}{p{0.22\linewidth}p{0.34\linewidth}p{0.34\linewidth}}
\toprule
question & Fixed SGQF & Zhao--Cui squared TT \\
\midrule
stored object & one Gaussian mean/covariance and an approximate innovation scalar & nonnegative TT approximation of filtering/posterior densities \\
main strength & deterministic fixed cloud, transparent scalar, analytical same-scalar gradient & richer non-Gaussian density representation and conditional maps \\
main limitation & Gaussian projection cannot store multimodality or skewed posterior shape & adaptive TT-cross/rank changes are not globally smooth; implementation is heavier \\
gradient lane & differentiates fixed Gaussian-projection likelihood & fixed-branch derivative differentiates a declared TT approximate scalar \\
best use & fast approximate likelihood, baseline, local diagnostic, candidate fixed target when branch diagnostics stay stable & high-dimensional non-Gaussian filtering and smoothing proposal \\
\bottomrule
\end{tabular}
\end{center}

The proposal is therefore not to replace Zhao--Cui with fixed SGQF.  The
proposal is to carry two honest candidates: a transparent fixed-rule Gaussian
projection likelihood with a same-scalar gradient, and a richer squared-TT
density approximation for the cases where Gaussian projection is too weak.

\section{Where Each Construction Comes From}
\label{sec:p31-source-map}

\begin{center}
\small
\begin{tabular}{p{0.34\linewidth}p{0.28\linewidth}p{0.28\linewidth}}
\toprule
construction & support & use in this note \\
\midrule
state-space and conceptual filtering equations & \citet[Section 2]{jia2012sparsegrid} & Sections~\ref{sec:p31-ssm}--\ref{sec:p31-gaussian-projection} \\
Gaussian approximation filter moment structure & \citet[Eq. 15--23]{jia2012sparsegrid} plus Proposition~\ref{prop:p31-affine-projection} & moment projection and update equations \\
one-dimensional and tensor-product Gaussian quadrature & \citet[Section 2.2.1]{jia2012sparsegrid} & Section~\ref{sec:p31-ghq} \\
sparse-grid formula and index band & \citet[Eq. 26--29]{jia2012sparsegrid} & Section~\ref{sec:p31-sgq-reconstruction} \\
point/weight generation and duplicate accumulation & \citet[Algorithm 1]{jia2012sparsegrid} & fixed cloud \(\mathcal C\) and toy grid \\
quadrature exactness and UKF relation & \citet[Theorems 3.1--3.2]{jia2012sparsegrid} & source-scope orientation, not posterior guarantee \\
adaptive sparse-grid mechanics & \citet[Section 3]{singh2018adaptive} & offline grid-design interpretation \\
fixed scalar and analytical gradient & BayesFilter project derivation & Sections~\ref{sec:p31-same-scalar}--\ref{sec:p31-gradient} \\
comparison with squared TT & Zhao--Cui companion note and \citet{zhao2024ttsequential} & Section~\ref{sec:p31-comparison} \\
\bottomrule
\end{tabular}
\end{center}

\section{Conclusion}
\label{sec:p31-conclusion}

FixedSGQF is now a well-defined standalone proposal: choose and save a
sparse-grid cloud, evaluate the Gaussian projection likelihood, differentiate
that same saved scalar, and reject branch changes that would make the value and
gradient disagree.  Its intellectual virtue is clarity.  Its mathematical
limitation is also clear: it is still a Gaussian projection filter.  That
honesty is what makes it a useful companion to the Zhao--Cui squared-TT lane
rather than a weaker substitute for it.

\bibliographystyle{plainnat}
\bibliography{../references}

\end{document}
